\section{Objectif}

L'utilisation d'algorithmes récursifs lors de l'écriture de programmes
\texttt{OCaml} est très répandue en raison de l'utilisation courante de
structures inductives, par exemple le type \texttt{list}. Malheureusement, ces
algorithmes font souvent face aux limitations imposées par la taille maximale de
la pile d'appels, contrairement aux fonctions implémentées avec des boucles. Par
exemple l'implémentation récursive de la fonction \texttt{sum} suivante est plus
limitée que l'implémentation utilisant des boucles :

\bigskip
\small
\begin{minipage}[t]{0.5\textwidth}
    \begin{minted}[breaklines=true]{ocaml}
(* Implémentation récursive *)
let rec sum (l : int list) : int =
    match l with
    | [] -> 0
    | x :: q -> x + sum q




(* Stack overflow during evaluation (looping recursion?). *)
sum (List.init 1_000_000 (fun i -> i))
    \end{minted}
\end{minipage}%
\begin{minipage}[t]{0.5\textwidth}
    \begin{minted}[breaklines=true]{ocaml}
(* Implémentation par boucles *)
let sum (l : int list) : int =
    let res = ref 0 in
    let l_ref = ref l in
    while !l_ref <> [] do
        res := !res + List.hd !l_ref;
        l_ref := List.tl !l_ref
    done;
    !res

(* 499999500000 *)
sum (List.init 1_000_000 (fun i -> i))
    \end{minted}
\end{minipage}
\normalsize

\bigskip

On s'intéresse alors à la conversion automatique de fonctions écrites de façon
récursive en des fonctions implémentées à l'aide de boucles. Il s'avère que la
conversion de fonctions récursives \textit{terminales} en fonctions implémentées
avec des boucles est assez simple, et cette optimisation est même déjà
implémentée dans les compilateurs \texttt{ocamlc} et \texttt{ocamlopt}. C'est
mon camarade Raphaël BOUÉ-MOREL qui a implémenté cette partie-là de la
conversion, et je me suis alors concentré sur la conversion de fonctions
récursives \textit{quelconques} en fonctions récursives \textit{terminales},
procédure appelée \textit{rectification} dans la suite.
